\documentclass[conference]{IEEEtran}
\usepackage{cite}
\usepackage{amsmath,amssymb}
\begin{document}
\title{Estimating Drink Temperature from Sound: A Smart Mug Prototype}
\author{\IEEEauthorblockN{Itsuki Morimoto}
\IEEEauthorblockA{\textit{Dept. of Electrical Engineering} \\
\textit{Kitamachi University}\\
Kitamachi, Japan \\
morimoto@example.org}
\and
\IEEEauthorblockN{Clara Westwood}
\IEEEauthorblockA{\textit{School of Fictional Engineering} \\
\textit{University of Nowhere}\\
Nowhere, Freedonia \\
westwood@example.org}}
\maketitle
\begin{abstract}
The pitch of the sound produced when a liquid is poured into a mug depends on its temperature through the viscosity and the speed of sound.
We present a fictional smart mug that estimates drink temperature from a short audio recording of pouring.
A lightweight convolutional network running on a microcontroller predicts the temperature with a mean absolute error of 2.3~\textdegree C across five drinks.
The method needs no contact sensor and can be added to existing mugs with a clip-on microphone.
\end{abstract}
\begin{IEEEkeywords}
acoustic sensing, embedded machine learning, smart objects
\end{IEEEkeywords}
\section{Introduction}
Smart drinkware promises to prevent burns and to keep drinks at a preferred temperature~\cite{morimoto2021,westwood2022}.
Most existing products measure temperature with a probe or a sensor in the base, which complicates cleaning and manufacturing.
Acoustic sensing offers a contact-free alternative, and it has been used to estimate the fill level of containers~\cite{morimoto2023}.

In this paper we show that the same principle can estimate temperature.
We collect a dataset of pouring sounds, train a compact model, and deploy it on a low-power microcontroller.

\section{Method}
Audio is recorded at 16~kHz.
\bibliographystyle{IEEEtran}
\bibliography{refs}
\end{document}
